% from ega0-21-fr.tex
\oldpage[0\textsubscript{IV-1}]{165}
  tout homomorphisme
  $u':A\to B'$ tel que $B'u'(\mathfrak m)=\mathfrak n'$ et tout
  $K$-isomorphisme $\sigma:L\to L'$ (relatif à $u_0$ et
  $u'_0=u'\otimes1:K\to L'$), il existe un isomorphisme $w:B\to B'$ tel que
  $u'=w\circ u$ et que $w_0=w\otimes1:L\to L'$ soit égal à $\sigma$.
  \item[\rm b$'$)] Pour tout homomorphisme $u':A\to B$ tel que
  $Bu'(\mathfrak m)=\mathfrak n$ et tel que $u'_0=u'\otimes1:K\to L$ soit égal
  à $u_0$, il existe un automorphisme $w$ de $B$ tel que $u'=w\circ u$ et que
  $w_0=w\otimes1:L\to L$ soit l'identité.
  \item[\rm c)] Désignant par $u_1:A/\mathfrak m^2\to B/\mathfrak n^2$
  l'homomorphisme déduit de $u$ par passage aux quotients, alors, pour tout
  homomorphisme local $u'_1:A/\mathfrak m^2\to B/\mathfrak n^2$ tel que
  $\operatorname{gr}_0(u'_1)=\operatorname{gr}_0(u_1)(=u_0)$, il existe un
  automorphisme $w_1$ de $B/\mathfrak n^2$ tel que
  $\operatorname{gr}_0(w_1)$ soit l'identité et que $u'_1=w_1\circ u_1$.
  \item[\rm c$'$)] Pour tout homomorphisme local $u'_1:A/\mathfrak m^2\to
  B/\mathfrak n^2$ tel que $\operatorname{gr}_0(u'_1)=\operatorname{gr}_0(u_1)$
  et $\operatorname{gr}_1(u'_1)=\operatorname{gr}_1(u_1)$
  (homomorphisme de $\mathfrak m/\mathfrak m^2$ dans
  $\mathfrak n/\mathfrak n^2$), il existe un automorphisme $w_1$ de
  $B/\mathfrak n^2$ tel que $\operatorname{gr}_0(w_1)$ soit l'identité et que
  $u'_1=w_1\circ u_1$.
\end{enumerate}
 Alors on a les implications $c)\Leftrightarrow c')\Leftrightarrow a)\Rightarrow b)
 \Rightarrow b')$. Si de plus $A$ est un anneau de Cohen, les cinq conditions
précédentes sont équivalentes.

\end{proposition}
\begin{proof}
Les implications $b)\Rightarrow b')$ et $c)\Rightarrow c')$ sont triviales.
Montrons que $a)$ implique $b)$.
L'homomorphisme $u$ fait de $B$ un $A$-module sans torsion, puisqu'il
transforme par hypothèse une uniformisante de $A$ en une uniformisante de $B$~;
il en résulte que $B$ est un $A$-module plat (0$_{\rm I}$, 6.3.4), donc une
$A$-algèbre de Cohen (19.8.1)~; le fait que $a)$ implique $b)$ est alors
 conséquence de (19.8.2, (i)) appliquée à $G=B'$, $\mathfrak J=\mathfrak n'$,
 $B'$ étant considéré comme $A$-algèbre pour $u'$ et l'homomorphisme
$B\to B'/\mathfrak n'$ étant le composé
$B\to B/\mathfrak n\to B'/\mathfrak n'$, qui est un $A$-homomorphisme en
 vertu de l'hypothèse $u'_0=\sigma\circ u_0$. Le même raisonnement prouve que
$a)$ entraîne $c)$, en prenant cette fois $G=B/\mathfrak n^2$,
$\mathfrak J=\mathfrak n/\mathfrak n^2$, $G$ étant considéré comme
$A$-algèbre pour $u'_1$.
Prouvons en second lieu que $c')$ implique $a)$. Les deux homomorphismes
$u_1$ et $u'_1$ sont tels que $u'_1=u_1+D$, où $D$ est une dérivation de
$A/\mathfrak m^2$ dans $\mathfrak n/\mathfrak n^2$ (20.1.1)~; d'ailleurs,
l'hypothèse $\operatorname{gr}_1(u'_1)=\operatorname{gr}_1(u_1)$ signifie
que $D$ est nulle dans $\mathfrak m/\mathfrak m^2$, et peut par suite être
considérée comme une dérivation de $K=A/\mathfrak m$ dans
$\mathfrak n/\mathfrak n^2$, et $\mathfrak n/\mathfrak n^2$ s'identifie
(par choix d'une uniformisante de $B$) au $K$-module $L$ (pour $u_0$).
D'autre part, les conditions imposées à $w_1$ entraînent que
 $\operatorname{gr}_1(w_1)$ est aussi l'identité (puisque
 $u_1(\mathfrak m/\mathfrak m^2)$ engendre
$\mathfrak n/\mathfrak n^2$)~; $w_1$ est donc de la forme
$x\mapsto x+D'(x)$, où $D'$ est cette fois une dérivation de
$B/\mathfrak n^2$ dans le $(B/\mathfrak n^2)$-module
$\mathfrak n/\mathfrak n^2$~; comme $w_1$ est l'identité dans
$\mathfrak n/\mathfrak n^2$, $D'$ peut encore (par l'identification
précédente de $\mathfrak n/\mathfrak n^2$ et de $L$) être considérée comme
une dérivation de $L$ dans $L$~; enfin, la relation $u'_1=w_1\circ u_1$
signifie que $D'$ prolonge $D$. Notons d'autre part que toute dérivation $D$
de $K$ dans $L$ correspond à un homomorphisme $u'_1$ vérifiant les conditions
de $c')$ (20.1.1)~; la condition $c')$ signifie donc que toute dérivation de
$K$ dans $L$ se prolonge en une dérivation de $L$ dans lui-même, c'est-à-dire
(20.6.5) que $L$ est séparable sur $K$.
Prouvons enfin que, lorsque $A$ est un anneau de Cohen, $b')$ entraîne $c')$.
Prouvons d'abord que sous les hypothèses de $c')$, il existe un homomorphisme
$u':A\to B$ qui vérifie les hypothèses de $b')$ et qui, par passage aux quotients,
donne $u'_1:A/\mathfrak m^2\to B/\mathfrak n^2$. En effet, cela résulte de
(19.8.6, (i)) appliqué à l'homomorphisme composé
$v':A\to A/\mathfrak m^2\xrightarrow{u'_1}B/\mathfrak n^2$~;
